This paper studies how housing markets price the physical depreciation of local public capital. Exploiting close-margin renewal-tax referendums on road maintenance in Ohio from 1991 through 2021, and a novel ConvNeXt-V2 vision model fine-tuned to measure physical road condition from satellite imagery, we document three findings.

First, the fiscal shock --- an 11-percent drop in maintenance spending --- is realized at the moment of the vote, but the response of housing prices is not. The pre-period RD coefficients are economically small and statistically indistinguishable from zero, ruling out anticipation. The contemporaneous and three-year-ahead coefficients are also indistinguishable from zero, ruling out rational-expectations capitalization. The price gap between narrowly-failing and narrowly-passing jurisdictions opens only in the fourth year following the vote, the same year in which physical road deterioration becomes statistically detectable in our computer-vision-based measurement. The long-run capitalized loss is approximately nine percent of the mean median sale price (about \$15{,}350 per house).

Second, the synchronization between the {\it visibility} of physical decay and the {\it timing} of price capitalization is decisive in distinguishing rational expectations from limited-attention and signal-extraction alternatives. We embed this distinction in a dynamic general-equilibrium model with a Bayesian belief-updating microfoundation, in which households observe the public-capital stock through a noisy signal whose precision shrinks endogenously with the magnitude of physical decay. The microfoundation nests rational expectations as a knife-edge case and identifies the friction parameter from a single transparent moment in the data. The calibrated belief-updating model matches both the dynamic price path and the cross-sectional heterogeneity by urban density and housing-price quantile without additional free parameters.

Third, the welfare implications are substantial. The capitalized housing-wealth loss exceeds the present value of the household's tax savings by roughly eight to one. The implied marginal value of public funds is $-1.4$, identifying the policy as welfare-destroying. The behavioral-welfare counterpart, evaluated under perfect foresight, places this number in the range $[-2.1, -1.6]$. The interpretation is that voters supported the levy cut not because they assessed it as net-beneficial under correct beliefs, but because the costs lay in a future state of the world that they did not observe at the moment of the vote.

\paragraph{Policy implications.}
Two policy levers are suggested by the analysis. First, mandatory pre-vote disclosure of road-condition assessments --- analogous to restaurant-grading systems, energy-efficiency disclosure at sale, and water-quality reporting --- would reduce the observability noise that drives the delayed-capitalization pattern documented here. The fire-protection findings of \citet{brasington2024fire} are directly relevant: a limited-awareness mechanism for capitalization of public services attenuates when awareness rises. Second, unbundling road-maintenance levies from other items on the ballot would reduce the cognitive cost emphasized in the limited-attention microfoundation of Section \ref{sec:fiscal_illusion}. The model-based counterfactual in Section \ref{sec:mvpf} suggests that even a modest reduction in the observability friction --- a doubling of pre-vote attention --- would reduce the steady-state capitalization gap by approximately forty percent.

\paragraph{Generality.}
The mechanism we document is plausibly common to many forms of slowly-depreciating public capital. Water mains, bridges, schools, and municipal IT systems share three features with the road-maintenance setting: their depreciation is gradual rather than discrete; their physical state becomes observable to residents and prospective buyers only with a substantial visibility lag; and the corresponding fiscal decisions occur in low-salience local elections. Where these features are present, our analysis predicts that markets will under-react to disinvestment relative to the rational-expectations benchmark, that welfare losses will accumulate in the form of a capitalization gap between wealth losses and tax savings, and that information-disclosure mechanisms targeting the observability lag will deliver welfare gains substantially exceeding their implementation cost.

\paragraph{Limitations and future work.}
Our analysis is conducted in a single state and a single class of public capital. Extending the framework to other geographies and other slowly-depreciating asset classes would test the generality of the visibility-lag mechanism. The belief-updating microfoundation in Section \ref{sec:belief_updating} treats observability noise as exogenous; an interesting extension would endogenize observability investment by local governments, voters, or third-party rating organizations. The computer-vision measurement framework introduced in Section \ref{sec:data} is itself extensible to other forms of physical public capital --- bridge condition, sidewalk integrity, urban tree cover --- where high-resolution aerial imagery is available. We leave these extensions to future work.
